% =====================================================================================================
% Section I. Scaling Laws and the Production Problem
% Writer: framework section writer, 2026-09-24. Sources: paper/notes/model_spec.md, paper/notes/theory_main_text.md,
% output/memos/m7_theory.md (+ review), paper/sections/appendix_proofs.tex, lit/SYNTHESIS.md section 4.
% Numbers: output/tables/m7_theory_ledger_checks.csv (Besiroglu parameters), appendix Lemma A1(v) (Kaplan sigma),
% output/tables/m7_theory_harberger.csv and m7 memo section 5 (Harberger accuracy). See paper/notes/writer_framework.md.
% =====================================================================================================

\section{Scaling Laws and the Production Problem}\label{sec:framework}

This section writes the scaling law as a production technology, sets out the cost side of a lab's training problem,
and translates the vocabulary of the scaling literature into that of production economics
(Table~\ref{tab:dictionary}). It then states three results used throughout: duality between the scaling law, the
compute-optimal allocation and the loss--compute frontier (Proposition~\ref{prop:duality}); gross complementarity as
a condition for interior compute optima (Lemma~\ref{lemma:sigma}); and revealed inference demand
(Proposition~\ref{prop:wedge}). Proofs are in Online Appendix~A.

\subsection{Technology, Cost, and Output}\label{sec:framework:tech}

A lab trains a language model with $N$ parameters on $D$ tokens (tokens processed, counting repetitions) and
measures its cross-entropy loss $L$, in nats per token, on a fixed held-out corpus. Following
\citet{hoffmann2022training}, and adding the productivity terms of the production-function literature, we write
\begin{equation}\label{eq:tech}
L=E+e^{-\omega}\Big[A\big(e^{\psi_N}N\big)^{-\alpha}+B\big(e^{\psi_D}D\big)^{-\beta}\Big]e^{\epsilon},
\qquad A,B,\alpha,\beta>0 .
\end{equation}
Here $E$ is the irreducible loss, and the power terms are the losses due to finite parameters and finite data.
Hicks-neutral productivity $\omega$ (an architecture or recipe that scales both terms) and parameter- and
data-augmenting productivity $\psi_N$ and $\psi_D$ (a cleaner corpus raises $\psi_D$) are known to the lab when it
chooses inputs: they are transmitted in the sense of \citet{marschak1944random}. Seed and evaluation noise $\epsilon$
is not \citep{zellner1966specification}. The Chinchilla law sets $\omega=\psi_N=\psi_D=\epsilon=0$.

Lower-case letters denote logs, $n\equiv\ln N$ and $d\equiv\ln D$. We call $R\equiv L-E$ the reducible loss and use
$y\equiv-\ln R$ as the log output index. Taking logs of \eqref{eq:tech} gives, exactly,
$y=\omega+F(n+\psi_N,d+\psi_D)-\epsilon$ with $F(n,d)\equiv-\ln(Ae^{-\alpha n}+Be^{-\beta d})$: a log production
function with Hicks-neutral productivity and an additive disturbance. With $u\equiv A(e^{\psi_N}N)^{-\alpha}$ and
$v\equiv B(e^{\psi_D}D)^{-\beta}$, the output elasticities are
\begin{equation}\label{eq:elast}
\varepsilon_N\equiv\frac{\partial y}{\partial n}=\frac{\alpha u}{u+v},\qquad
\varepsilon_D\equiv\frac{\partial y}{\partial d}=\frac{\beta v}{u+v},
\end{equation}
and neither depends on $\omega$, $\epsilon$ or $E$. When $\alpha=\beta$, $R^{-1/\alpha}$ is a constant-returns CES
aggregate of $N$ and $D$. When $\alpha\neq\beta$, the technology is directly additive and non-homothetic: the
elasticity of substitution varies with the input mix, and the compute-optimal ratio of tokens to parameters drifts
with scale. We also use the generalized family that raises the bracket in \eqref{eq:tech} to a power $\kappa>0$,
with inner exponents $\alpha$ and $\beta$. Chinchilla is $\kappa=1$; Kaplan's law is
$\kappa\approx0.103$ with inner exponents $0.738$ and $1$ and $E=0$. The outer exponent scales both elasticities by
$\kappa$ and leaves the isoquants unchanged, so it separates returns to scale from substitution; $\kappa=1$ is a
testable restriction (Section~\ref{sec:tech}).

Training a dense transformer takes about six floating-point operations (FLOPs) per parameter per token, so training
compute is $C=6ND$ and $c\equiv\ln C=\ln6+n+d$ \citep{kaplan2020scaling}; Section~\ref{sec:tech} treats the
attention and embedding FLOPs this approximation omits. The cost is multiplicative rather than linear: both log-cost
weights, $\partial c/\partial n$ and $\partial c/\partial d$, equal one, and the implicit price of a parameter, $6D$,
rises with the quantity of the other input. Two consequences follow. First, minimizing compute for a given loss
equates the output elasticities, $\varepsilon_N=\varepsilon_D$. The first-order conditions behind the share-equation
approaches of \citet{deloecker2012markups} and \citet{gandhi2020identification} link output elasticities to cost
shares; here the shares are known from FLOP accounting, so no price data are needed. Second, isocosts are straight
lines in logs, $n+d=c-\ln6$, which sharpens the curvature condition for an interior optimum
(Section~\ref{sec:framework:sigma}).

A deployed model spends about $2N$ FLOPs per token it processes, so a lab that expects to serve $T$ tokens over the
model's life faces lifetime compute $6ND+2NT$ \citep{sardana2024beyond}. We take lifetime compute as the lab's
cost, minimized subject to a loss target $L\le\bar L$ and possibly a data constraint $D\le\bar D$ or a memory
constraint $N\le\bar N$. In this language, iso-loss curves are isoquants, IsoFLOP profiles are isocosts, the
compute-optimal allocation is a conditional factor demand, and the loss--compute frontier is the inverse of a cost
function (Figure~\ref{fig:geometry}, panel (a)). We write $M\equiv D/N$ for tokens per parameter and $M^*(C)$ for the
training-optimal ratio at compute $C$.

Loss is measured in units specific to a tokenizer and an evaluation corpus, so it matters which objects depend on
them. Isoquants, the elasticity of substitution, the compute-optimal allocation and the wedge defined below are
ordinal: they are unchanged by any strictly increasing transformation of output (loss, bits per byte or a monotone
benchmark link) and do not depend on $E$, $\omega$ or $\kappa$ (Online Appendix~A, Lemma~\ref{lem:sigma};
Proposition~\ref{prop:wedge} below). The
frontier elasticity, returns to scale and productivity levels are cardinal: they are defined in units of reducible
loss. Input units matter too: $M^*$ counts tokens and so depends on the tokenizer, while the exponents do not.

\subsection{A Dictionary}\label{sec:framework:dict}

Table~\ref{tab:dictionary} lists eighteen correspondences between scaling-law concepts and concepts of production
theory and IO econometrics. The cost-side entries in panel A are exact: once $C=6ND$, they are identities of duality
theory. The others are close analogies: the econometric problem is the same, but its source can differ. The
collinearity of inputs along the expansion path, for example, is generated by optimizing behavior in cross-lab data
and by design choices in single-lab sweeps (Section~\ref{sec:ident}).

Three analogies break down. There are no market prices for parameters and tokens; returns to scale measured in loss
depend on the cardinalization of output, so ``diminishing returns to compute'' has no unit-free content; and the
inputs of a designed sweep are not chosen by optimizing firms, so the endogeneity machinery of IO applies to cross-lab
data (Section~\ref{sec:obs}), not to single-lab experiments. Many entries have precursors in machine learning,
reviewed in the introduction; the table links each to the econometric problem it instantiates and to that
literature's remedies.

% Table 1: compact dictionary between scaling-law objects and production/IO econometrics (v2, <= 8 rows).
% Writer: framework section writer (v2), 2026-09-24. Replaces the 18-row v1 table (sections/v1/framework.tex).
% Corrections after referee round 1 (R1 comment 8; R2 Major 9a, 9d; R3 M1, minor 10):
%   - on-path collinearity is Marschak-Andrews / Bond-Soderbom collinearity under common prices, not ACF functional
%     dependence (R1 8a); IsoFLOP sweeps are assigned (designed) inputs, not price shocks (R1 8d; R2 9d);
%   - no GNR label for the flagship normalization; no TFPR row for bounded benchmark links (R1 8e, 8h);
%   - Nerlove (1963) is cited for the returns-to-scale (cost-elasticity) measure only, not for mismeasurement;
%     scale-correlated input mismeasurement is attributed to Basu-Fernald (R1 8g);
%   - "Approach 1" renamed: the frontier through IsoFLOP minima is the "IsoFLOP-minima frontier" (R2 9a);
%   - the wedge is Raval's ratio statistic with FLOP-accounting weights (R1 3a).
\begin{table}[tp]
\centering
\caption{Scaling Laws in the Language of Production Economics}
\label{tab:dictionary}
\footnotesize
\setlength{\tabcolsep}{3pt}
\begin{tabular}{@{}>{\raggedright\arraybackslash}p{0.22\textwidth}>{\raggedright\arraybackslash}p{0.22\textwidth}>{\raggedright\arraybackslash}p{0.325\textwidth}>{\raggedright\arraybackslash}p{0.185\textwidth}@{}}
\toprule
Scaling-law object & Production or IO counterpart & Use in this paper & References \\
\midrule
Iso-loss curve; IsoFLOP profile ($6ND=C$) & Isoquant; output along an isocost with unit log-cost weights & The optimum equates output elasticities without price data; interior optima need $\sigma<1$ (Lemma~\ref{lemma:sigma}) & \citealt{shephard1953cost} \\[3pt]
Chinchilla law and its $\kappa$ family & Quasi-homothetic technology; outer exponent $\kappa$ sets returns to scale & All curvature lies across the expansion path; $\kappa=1$ is testable (Section~\ref{sec:tech}) & \citealt{hoffmann2022training}; \citealt{kaplan2020scaling} \\[3pt]
Compute-optimal $N^*(C)$, $D^*(C)$; IsoFLOP-minima frontier & Conditional factor demands; inverse cost function & Duality links $a$, $\gamma$ and $\sigma^*$; the elasticity of cost with respect to reducible loss is $-1/\gamma$ (Proposition~\ref{prop:duality}) & \citealt{uzawa1964duality}; \citealt{nerlove1963returns} \\[3pt]
Compute-optimal ladders; fixed tokens-per-parameter designs & Collinear inputs under common prices & Identify the path and the frontier, not $\sigma^*$; information about curvature is fourth order in allocation errors (Proposition~\ref{prop:ident}) & \citealt{marschak1944random}; \citealt{bond2005adjustment} \\[3pt]
IsoFLOP sweeps; factorial designs & Assigned variation in the input mix; optimal design & Curvature at the IsoFLOP minimum identifies $\sigma^*$ without a functional form (Proposition~\ref{prop:modelfree}) & \citealt{box1959design}; \citealt{kiefer1959optimum} \\[3pt]
Over-training, $M>M^*(C)$ & Input-specific wedge: ratio of output to cost elasticities & Reveals the value of compactness and its expenditure share $s$ (Proposition~\ref{prop:wedge}; Section~\ref{sec:wedge}) & \citealt{deloecker2012markups}; \citealt{raval2023testing} \\[3pt]
Lab-specific recipes, beliefs and output concepts & Factor-augmenting heterogeneity; output-concept mismatch & Contaminate the wedge: $\hat w=w\,w_L/w_{L'}$ (Proposition~\ref{prop:wedge}) & \citealt{demirer2020production}; \citealt{bond2020unpleasant} \\[3pt]
Embedding counts; omitted FLOPs; untuned hyperparameters & Scale-correlated input mismeasurement; flexible inputs not re-optimized & Shift the allocation exponent $a$ and $M^*$ (Section~\ref{sec:tech}) & \citealt{basu1997returns}; \citealt{porian2024resolving} \\
\bottomrule
\end{tabular}
\par\smallskip
\begin{tablenotes}
The first three rows are exact once training compute is $C=6ND$; the others are analogies in which the econometric
problem is the same but its source can differ. $M=D/N$ is tokens per parameter and $M^*(C)$ its compute-optimal value;
$w=\varepsilon_N/\varepsilon_D$ is the wedge and $s=(w-1)/w$; $a$ is the allocation exponent, $\gamma$ the frontier
elasticity, $\sigma^*$ the elasticity of substitution on the expansion path and $\kappa$ the outer exponent; $L$ and
$L'$ are the technologies used by the econometrician and by the developer.
\end{tablenotes}
\begin{tablenotes}[Source]
Authors' compilation.
\end{tablenotes}
\end{table}


\subsection{Duality}\label{sec:framework:duality}

\begin{proposition}[Duality]\label{prop:duality}
Let $\kappa=1$ and $\psi_N=\psi_D=0$ in \eqref{eq:tech}.
\begin{enumerate}
\item[(i)] The allocation that minimizes training compute $6ND$ subject to $L\le\bar L$, or equivalently maximizes
$y$ at given $C$, satisfies $\varepsilon_N=\varepsilon_D$. The conditional factor demands are
\[
N^*(C)=G\,(C/6)^{a},\qquad D^*(C)=G^{-1}(C/6)^{b},
\]
with $a\equiv\beta/(\alpha+\beta)$, $b\equiv1-a$ and $G\equiv(\alpha A/\beta B)^{1/(\alpha+\beta)}$, so that
$M^*(C)=G^{-2}(C/6)^{b-a}$.
\item[(ii)] The loss--compute frontier and the cost function are
\[
L^*(C)=E+K\,(C/6)^{-\gamma},\qquad C^*(L)=6\,\big[K/(L-E)\big]^{1/\gamma},
\]
with $\gamma\equiv\alpha\beta/(\alpha+\beta)$ and $K\equiv\frac{\alpha+\beta}{\beta}AG^{-\alpha}$. Along the
expansion path, $\varepsilon_N=\varepsilon_D=\gamma$.
\item[(iii)] With productivity, $y^*=\omega+\gamma(c-\ln6+\psi_N+\psi_D)-\ln K$ and
\[
\ln M^*(C)=-2\ln G+(b-a)(c-\ln6)-\frac{2\chi}{\alpha+\beta},\qquad \chi\equiv\beta\psi_D-\alpha\psi_N .
\]
At given compute, Hicks-neutral productivity $\omega$ does not affect the allocation; better data (higher
$\psi_D$) raise $N^*$ and lower both $D^*$ and $M^*$.
\end{enumerate}
\end{proposition}

\noindent The proof is in Online Appendix~A (Lemma~\ref{lem:alloc}). The proposition is the scaling-law instance of
Shephard--Uzawa duality \citep{shephard1953cost,uzawa1964duality}: the technology, the conditional factor demands and
the cost function carry the same information, linked by $a=\beta/(\alpha+\beta)$ and
$\gamma=\alpha\beta/(\alpha+\beta)$. Allowing $\kappa\neq1$ leaves part (i) unchanged and multiplies $\gamma$ by
$\kappa$. By part (ii), $-1/\gamma$ is the elasticity of compute with respect to reducible loss along the frontier,
the object \citet{nerlove1963returns} estimated for electricity supply. Under the parameters of
\citet{besiroglu2024chinchilla}, $\gamma=0.178$: a one percent reduction in reducible loss requires about 5.6 percent
more compute.

The proposition also classifies the three estimation strategies of \citet{hoffmann2022training}. Approach~1, the
lower envelope of loss against compute across runs of fixed size, estimates the cost function in part (ii).
Approach~2, the regression of the IsoFLOP-minimizing $\ln N$ on $\ln C$, estimates the conditional factor demand in
part (i), which identifies $a$ and $G$ but not $\alpha$ and $\beta$. Approach~3, a fit of \eqref{eq:tech} to all
runs, is a primal estimator. Together they form a system of a production function and its first-order conditions
\citep{klump2007factor,leonledesma2010identifying}, and their differences are a test of duality.
\citet{hoffmann2022training} report values of $a$ between 0.46 and 0.50 across the approaches;
Section~\ref{sec:tech} asks whether such gaps exceed sampling error.

Part (iii) holds at given compute. Because Hicks-neutral productivity does not tilt the expansion path, the input mix
cannot proxy for it, and the control-function inversion of \citet{olley1996dynamics} is unavailable for the
productivity that differs most across labs (Section~\ref{sec:ident}). At a given loss target, by contrast, better data
leave $N^*$ unchanged and lower $D^*$ one for one, and Hicks-neutral productivity moves the chosen mix when
$\alpha\neq\beta$, because the technology is non-homothetic (Online Appendix~A, Lemma~\ref{lem:alloc}(v)).

\subsection{Substitution}\label{sec:framework:sigma}

For the technology \eqref{eq:tech}, the Hicks elasticity of substitution between parameters and data is
\begin{equation}\label{eq:sigma}
\sigma=\frac{\alpha u+\beta v}{\alpha u(1+\beta)+\beta v(1+\alpha)},
\end{equation}
a formula derived by \citet{hao2026theory}. Because $u$ and $v$ enter only through their ratio, $\sigma$ depends on
the input mix only through the wedge $w\equiv\varepsilon_N/\varepsilon_D=\alpha u/(\beta v)$:
\begin{equation}\label{eq:sigmaw}
\sigma(w)=\frac{1+w}{1+\alpha+w(1+\beta)},\qquad
\sigma^*\equiv\sigma(1)=\frac{2}{2+\alpha+\beta}.
\end{equation}
The local elasticity lies between $1/(1+\alpha)$ and $1/(1+\beta)$ and equals $\sigma^*$ on the compute-optimal path,
where $w=1$. Under the \citet{besiroglu2024chinchilla} parameters ($\alpha=0.348$, $\beta=0.366$), $\sigma^*=0.737$,
and $\sigma$ varies only between 0.732 and 0.742 over all input mixes. Because $\sigma^*=1/[1+(\alpha+\beta)/2]$, the
value of 0.762 that \citet{hao2026theory} obtain by averaging the exponents of \citet{hoffmann2022training} is exactly
$\sigma^*$ for those exponents. The joint law of \citet{kaplan2020scaling} implies $\sigma\in[0.500,0.575]$ and
$\sigma^*=0.535$ (Online Appendix~A, Lemma~\ref{lem:sigma}): the two canonical laws disagree about the curvature of
the isoquants, not only about the expansion path.

\begin{lemma}[Compute-optimal training reveals gross complementarity]\label{lemma:sigma}
Let $y=F(n,d)$ be any twice continuously differentiable technology with $F_n,F_d>0$, and let training compute be
$C=6ND$. At an allocation with $\varepsilon_N=\varepsilon_D$:
\begin{enumerate}
\item[(i)] the second-order condition for minimizing compute at given output holds strictly if and only if
$0<\sigma<1$ there;
\item[(ii)] if $\sigma>1$ or $\sigma<0$ there, the allocation minimizes rather than maximizes output along its
isocost, and if $\sigma>1$ everywhere, compute-optimal training is a corner solution.
\end{enumerate}
Hence every interior compute-optimal allocation at which the second-order condition holds strictly has $\sigma<1$.
In the generalized family, $\sigma<1$ at every input mix, so the first-order conditions of
Propositions~\ref{prop:duality} and~\ref{prop:wedge} are sufficient.
\end{lemma}

\noindent The proof is in Online Appendix~A (Lemma~\ref{lem:soc} and Corollary~\ref{cor:complements}). With a linear
cost $p_NN+p_DD$, isocosts are straight lines in levels, and any convex isoquant ($\sigma>0$) gives an interior
tangency. With $C=6ND$ they are straight lines in logs, so the isoquant must be convex in logs, which requires
$\sigma<1$. Economically, the implicit price of each input rises with the quantity of the other, so shifting toward
one input raises its own price; that pays only if the inputs are complements. The existence of a compute-optimal
frontier, on which neither very large under-trained nor very small over-trained models are efficient, is therefore
evidence of gross complementarity. The families we fit impose $\sigma<1$, but on the
best-designed sweep, estimates from forms that do not impose it also imply $\sigma<1$ (Section~\ref{sec:tech}).

The lemma has two further uses. Observational-equivalence arguments can move $\sigma^*$ anywhere in $(0,1)$ but
never to one or above, which bounds the non-identification results of Section~\ref{sec:ident}. And $\sigma>0$ is what
makes over-training a sensible response to inference demand: under the Leontief approximation of
\citet{hao2026theory} ($\sigma=0$), the optimal allocation sits at the kink whatever the inference load, whereas under
complementarity trading parameters for tokens is feasible at a cost governed by $\sigma^*$ (Figure~\ref{fig:geometry},
panel (b)).

\begin{figure}[tp]
\centering
\includegraphics[width=\textwidth]{m7_theory_geometry}
\caption{The Geometry of the Training Problem}\label{fig:geometry}
\begin{figurenotes}
Panel (a): the $(\log_{10}N,\log_{10}D)$ plane under the \citet{besiroglu2024chinchilla} parameters. Solid gray curves
are iso-loss curves (isoquants) at the frontier losses for $C=10^{20},\dots,10^{24}$ FLOP; dotted lines are isocosts
$6ND=C$; the blue line is the compute-optimal expansion path. The orange dot is a run with $N=8\times10^{9}$ and
$D=1.5\times10^{13}$ (the configuration of Llama~3~8B), joined to the path at equal compute; its wedge is $w=5.2$.
Panel (b): three members of the generalized family, with $\sigma^*=0.60$, $0.74$ and $0.85$, constructed to share the
expansion path and the loss--compute frontier (Online Appendix~A, Proposition~\ref{prop:fd}(iv)). The vertical axis
is reducible loss in excess of the frontier at $C=10^{22}$ FLOP, as a function of the deviation of tokens per
parameter from the compute-optimal ratio. The technologies coincide on the path and differ off it.
\end{figurenotes}
\end{figure}

\subsection{Wedges, Inference Demand, and Efficiency}\label{sec:framework:wedge}

\begin{proposition}[Revealed inference demand]\label{prop:wedge}
A lab chooses $(N,D)$ to minimize lifetime compute $6ND+2NT$ subject to $L\le\bar L$, where $L$ belongs to the
generalized family with inner exponents $\alpha$ and $\beta$.
\begin{enumerate}
\item[(i)] If neither a data nor a memory constraint binds, the optimum satisfies
\[
w\equiv\frac{\varepsilon_N}{\varepsilon_D}=1+\frac{T}{3D}=\frac{6ND+2NT}{6ND},
\]
so the wedge equals the ratio of lifetime to training compute, and $T=3D(w-1)$. The wedge is invariant to
$\omega$, $\epsilon$, $E$ and $\kappa$.
\item[(ii)] For a model with compute $C$ and ratio $M$,
\[
\ln w=\Big(\frac{1}{\sigma^*}-1\Big)\ln\frac{M}{M^*(C)},\qquad\text{so}\qquad
\frac TD=3\bigg[\Big(\frac{M}{M^*(C)}\Big)^{1/\sigma^*-1}-1\bigg],
\]
where $M^*(C)$ is the training-optimal ratio at the model's own compute.
\item[(iii)] Relative to the training-only frontier, the model uses
\[
\frac{C}{C_{\min}}=\Big(\frac{\alpha+\beta w}{\alpha+\beta}\Big)^{(\alpha+\beta)/(\alpha\beta)}w^{-1/\alpha}
\]
times the least training compute that attains its loss. This ratio depends on $(w,\alpha,\beta)$ alone, not on
$A$, $B$, $E$ or productivity; with $\kappa=1$ the first exponent is $1/\gamma$.
\item[(iv)] A binding data constraint gives $w=1/(\theta_D+\mu)<1+T/(3D)$, where $\theta_D\equiv6D/(6D+2T)$ and
$\mu>0$ is the constraint's shadow cost, so that $w<1$ is possible; a binding memory constraint gives
$w=(1+\nu)(1+T/3D)>1+T/(3D)$ with shadow cost $\nu>0$. If labs differ in factor-augmenting productivity, the wedge
computed with a common technology is $\hat w=we^{-\chi}$.
\end{enumerate}
\end{proposition}

\noindent The proof is in Online Appendix~A (Proposition~\ref{prop:A-wedge}, Lemmas~\ref{lem:geometry}
and~\ref{lem:ce}, and Corollary~\ref{cor:suff}). Each parameter is used twice: in training, where it costs $6D$ FLOPs, and in serving, where it
costs $2T$ more; tokens are used only in training. At the lifetime optimum the lab equates the ratio of output
elasticities to the ratio of lifetime log-cost weights, $1:\theta_D$. An econometrician who sees only training cost,
with weights $1:1$, observes a wedge: the lab acts as if parameters were $w$ times as expensive, relative to tokens,
as training accounting says. Inverting the first-order condition makes the wedge a revealed quantity: $T/D=3(w-1)$ is
anticipated lifetime inference per training token, and $w-1$ is anticipated lifetime inference compute relative to
training compute. \citet{sardana2024beyond} solve the forward problem of choosing $N$ and $D$ given $T$;
Proposition~\ref{prop:wedge} runs it backward and needs neither the lab's productivity nor its loss.

The construction is the production approach of \citet{deloecker2012markups} applied to a different wedge. They
recover a markup as the ratio of a variable input's output elasticity to its revenue share, using cost minimization
and no demand model; here the ratio of two output elasticities, set against the ratio of their training-cost weights,
recovers the cost of a second use of one input. The wedge is input-specific, like the capital wedge of
\citet{hsieh2009misallocation}, but it is neither a distortion nor market power, and it is not a markup: it is the
shadow cost of serving, which optimal behavior should reveal. The analogy carries over one known weakness of the
production approach \citep{raval2023testing,demirer2020production}: by part (iv), factor-biased productivity
contaminates the wedge, as labor-augmenting technology contaminates markups estimated from labor.

By part (ii), revealed inference demand depends on the technology only through $M^*(C)$, the training-optimal ratio
at the model's own compute, and $\sigma^*$. The log wedge is the model's
distance from the expansion path at equal compute (Figure~\ref{fig:geometry}, panel (a)), scaled by curvature. A model
with 8 billion parameters trained on 15 trillion tokens, the configuration of Llama~3~8B, has $M\approx1{,}875$ and,
under the \citet{besiroglu2024chinchilla} parameters, $w\approx5.2$ (5.3 under the reference technology of
Section~\ref{sec:wedge}): it is built as if the lab expected about 12.7 inference tokens per training token, roughly $1.9\times10^{14}$ tokens over its life. Section~\ref{sec:wedge}
estimates such wedges, with confidence bands, for 173 open-weight models; their uncertainty is dominated by $M^*(C)$,
which on-path data identify only by assuming optimality (Section~\ref{sec:ident}).

Part (iii) is \citeauthor{farrell1957measurement}'s \citeyearpar{farrell1957measurement} input-oriented measure of
allocative efficiency relative to the training-only frontier. The model above uses 5.5 times the training compute
that a compute-optimal allocation would need for its loss, a Farrell cost efficiency of 0.18. That is not waste in
lifetime terms but the training-cost price of a model that is cheaper to serve. To second order,
$\ln(C/C_{\min})\approx\sigma^*(\ln w)^2/[4(1-\sigma^*)]$, a Harberger triangle whose curvature is set by $\sigma^*$.
The approximation is accurate to within 1 percent for $|\ln w|\le0.7$ but overstates the loss by 23 percent at
$w=5.2$, so we use the exact formula.

Wedges below one cannot come from inference demand. Gopher (280 billion parameters, 300 billion tokens) has $w=0.36$
under the same parameters and used 2.0 times the compute its loss required; GPT-3 has $w=0.43$. Such wedges reveal
a binding data constraint (part iv), data-augmenting productivity, or a belief error, here optimization against the
\citet{kaplan2020scaling} law; Section~\ref{sec:obs} separates this allocative component from technical change. More
generally, the rival sources of a wedge have known signs: data scarcity and data-augmenting productivity ($\chi>0$)
lower the measured wedge and so understate $T$, while memory or latency constraints and parameter-augmenting
productivity raise it. Section~\ref{sec:wedge} evaluates each.
